\section{Reparametrization to arclength}\label{sec:arclength-reparam}

\begin{theorem}\label[theorem]{thm:arclength}
Let $s\in \intervaloc{1/2,1}$. The (rescaling and) reparametrization operator 
\begin{align*}
\ARC: H^{1+s}_\reg\left( \R/\Z,\R^n \right) \to H^{1+s}_\unitspeed\left( \R/\Z,\R^n \right),
\quad	\gamma &\mapsto \frac{1}{\mathcal{L}(\gamma)} \, \gamma\circ\psi^{-1}_\gamma,
\end{align*} 
from the regular $H^{1+s}$-curves to the arclength
parametrized $H^{1+s}$-curves 
is continuous with respect to the $H^{1+s}$ topology, where 
$\psi_{\gamma}(t) \ceq \frac{1}{\mathcal{L}(\gamma)}\mathcal{L}(\gamma\mid_{[0,t]})$ for $t\in [0,1]$.
	\end{theorem}
Similar results have been shown in different settings; 
see, e.g., \cite[Theorem 2]{anosov_1981}. We follow the 
approach of Reiter in \cite[Theorem 1.13]{reiter_2009} who proved this result in the case $s=1$.
	Note that the result stated also holds for non-injective curves.

We prove \Cref{thm:arclength} by showing that the mapping $\ARC$ 
is the composition of 
continuous functions, dealt with in 
three separate lemmas. For that we define
\begin{align*}
	H^{1+s}_\homeom((0,1)) 
	&\ceq 
	\braces[\big]{
		\psi\in H^{1+s}((0,1)) 
		:  
		\psi'(x)>0 \Foa x\in[0,1],\; \psi(0)=0,\; \psi(1)=1
	}
	,
	\\
	H^{1+s}_\periodic((0,1)) 
	&\ceq 
	\braces[\big]{
		\psi\in H^{1+s}_\homeom((0,1)) 
		:  
		\psi'(0)=\psi'(1)
	}.
	\end{align*}
	The space $H^{1+s}_\homeom((0,1))$ is the set of orientation-preserving 
	$H^{1+s}$-diffeomorphisms of  the interval $[0,1]$, and its
	subspace $H^{1+s}_\periodic((0,1))$ consists of those diffeomorphisms, that 
	have a $1$-periodic derivative.
	The first step is to show continuity of the inversion operator: 
	\begin{lemma}
		\label[lemma]{inversioniscontinuous}
		For any $s\in (\frac{1}{2},1]$ the inversion operator $^{-1}:H^{1+s}_\homeom((0,1))\to H^{1+s}_\homeom((0,1))$ 
		is well-defined and continuous. In addition, we have for any 
		$\psi\in H^{1+s}_\homeom((0,1))$, $(\psi^{-1})'=\frac{1}{\psi'\circ\psi^{-1}}$.
	\end{lemma}
	\begin{proof}
For $s=1$ this is proven in 
\cite[Theorem 1.16]{reiter_2009}, so we may assume that $s\in (\frac{1}{2},1)$.
Since $H^{1+s}\equiv H^{1+s}((0,1))$ 
embeds continuously into $C^1([0,1])$, any 
$\psi\in H^{1+s}_e\equiv H^{1+s}_e((0,1))$ is a diffeomorphism
with a  continuously differentiable inverse $\phi \ceq \psi^{-1}$
satisfying $\phi(0)=0$, $\phi(1)=1$, and $\phi'=\frac{1}{\psi'\circ\phi}\ge 
\frac{1}{\left\|\psi' \right\|_{C^0}}>0$. 
By means of the change of variables $x=\psi(w),$ $y=\psi(v)$ with $
\d x=\psi'(w)\d w$ and $\d y=\psi'(v)\d v$ we can express the seminorm
of $\psi'\circ\phi$ as
\begin{align}
\lfloor\psi' &\circ\phi\rfloor_{s,2}^2
\textstyle
=\iint_{[0,1]^2}\frac{\left|\psi'\circ \phi(x) - \psi'\circ \phi (y) \right|^2}{\left|x-y \right|^{1+2s}}\d x\d y 
= \iint_{[0,1]^2}\frac{\left|\psi'(w) - \psi'(v) \right|^2}{\left|\psi(w)-\psi(v) \right|^{1+2s}}\psi'(w)\psi'(u)\d w\d v\notag\\
&
\textstyle
\le  \left\|\psi' \right\|^2_{C^0}\iint_{[0,1]^2}\frac{\left|\psi'(w) 
- \psi'(v) \right|^2}{\left|\psi(w)-\psi(v) \right|^{1+2s}}\d w\d v\le
c_\psi^{-(1+2s)}\left\|\psi' \right\|^2_{C^0}
\lfloor\psi'\rfloor_{s,2}^2<\infty,\label{eq:erste-chain}
\end{align}
where we used the bilipschitz estimate $|\psi(w)-\psi(v)|\ge c_\psi|w-v|$ 
for
all $v,w\in [0,1]$ for the diffeomorphism $\psi\in H^{1+s}_e$.
Since the map $r:[c,\infty)\to(0,\frac{1}{c}],$ $r(t)=\frac{1}{t}$ 
is Lipschitz for any $c>0$, the derivative $\phi'=r\circ\psi'\circ
\phi$ is of class $H^{s}((0,1))$ and therefore $\phi\in H^{1+s}$, so that
the inversion operator is well-defined. 

Now consider a sequence $(\psi_k)_k\subset H^{1+s}_\homeom$ converging 
to $\psi_\infty$ in $H^{1+s}_\homeom$ as $k\to\infty$. 
By virtue of the compact embedding $H^{1+s}\hookrightarrow C^1$ we
may assume without loss of generality 
that 
there is some constant $c>0$ such that $\psi'_k(x)\ge c$ for all $k\in\N$.
We want to prove that  $\phi_k \ceq \psi^{-1}_{k}$ converge to $\phi_\infty \ceq 
\psi^{-1}_\infty$ in $H^{1+s}$ as $k\to\infty$.
With 
\begin{align*}
	\left\|\phi_k-\phi_\infty \right\|_{C^0}
	&=\left\|\phi_\infty\circ\psi_\infty\circ\phi_k-\phi_\infty
	\circ\psi_k\circ\phi_k \right\|_{C^0}
	=\left\|\phi_\infty\circ\psi_\infty-\phi_\infty\psi_k \right\|_{C^0},
\end{align*}
where we used the fact that $\phi_k$ is a homeomorphism of 
$[0,1]$ onto itself, we obtain $
\left\|\phi_k-\phi_\infty \right\|_{C^0}
\le \left\|\phi_\infty' \right\|_{C^0}
\left\|\psi_k-\psi_\infty \right\|_{C^0}\to 0$ as $k\to\infty$.
		Furthermore,  
		\begin{align*}
\left\|\phi'_k-\phi_\infty' \right\|_{C^0}&\textstyle
\le \frac1{c^2}\left\|\psi'_k\circ\phi_k-\psi_\infty'
\circ \phi_\infty \right\|_{C^0}\\
			&
			\textstyle
\le \frac1{c^2}\left( \left\|\psi'_k\circ\phi_k-\psi_\infty'
\circ \phi_k \right\|_{C^0}+\left\|\psi_\infty'
\circ\phi_k-\psi_\infty'\circ \phi_\infty \right\| \right),
		\end{align*}
where the first summand on the right-hand side
equals $\left\|\psi'_k-\psi_\infty' \right\|_{C^0}\to 0$, 
since $\phi_k$ is a diffeomorphism and $\psi_k\to\psi_\infty$ 
in $C^1([0,1])$ as $k\to\infty$. 
The second summand goes to zero, as $\psi_\infty'$ is uniformly 
continuous on $[0,1]$, 
and we already showed that $\phi_k\to\phi_\infty$ 
in $C^0([0,1])$ as $k\to\infty$. 
It remains to show convergence with respect to the fractional Sobolev norm. 
\begin{align}\label{eq:seminorm-basic-estimate}
&\textstyle
\lfloor\phi'_k-\phi'_\infty\rfloor_{s,2}
\textstyle
=  \big\lfloor\frac{\psi'_k\circ\phi_k-\psi_\infty'\circ 
\phi_\infty}{(\psi'_k\circ\phi_k)\cdot 
(\psi_\infty'\circ\phi_\infty)}\big\rfloor_{s,2}\\
			&\le
			\textstyle
\lfloor\psi'_k\circ\phi_k\!-\!\psi_\infty'
\circ \phi_\infty\rfloor_{s,2}
\big\|\frac{1}{(\psi'_k\circ \phi_k)\cdot (\psi_\infty'\circ\phi_\infty)} 
			\big\|_{C^0}
+ \|\psi'_k\circ\phi_k\!-\!\psi_\infty'\circ \phi_\infty\|_{C^0}
\big\lfloor\frac{1}{(\psi'_k\circ\phi_k)
\cdot (\psi_\infty'\circ\phi_\infty)}\big\rfloor_{s,2},\notag
		\end{align}
where we used \cite[Proposition A.5]{knappmann-etal_2023}.
The second summand on the right-hand side of 
\cref{eq:seminorm-basic-estimate} converges to zero as $k\to\infty$, 
since the $C^0$-norm equals $\|\psi_k'-\psi_\infty'\|_{C^0}\to 0$, and
the seminorm is bounded from above by 
\begin{align}\label{eq:upper-bound-second-seminorm}
\textstyle
\frac1{c^2}\lfloor(\psi_k'\circ\phi_k)
\cdot (\psi_\infty'\circ\phi_\infty)\rfloor_{s,2}&
\textstyle
\le\frac1{c^2}\big(
\|\psi_k'\|_{C^0}\lfloor\psi_\infty'\circ\phi_\infty\rfloor_{s,2}+
\lfloor\psi_k'\circ\phi_k\rfloor_{s,2}\|\psi_\infty'\|_{C^0}\big)\\
&
\textstyle
\overset{\cref{eq:erste-chain}}{\le}
\frac1{c^2}\|\psi_k'\|_{C^0}\|\psi_\infty'\|_{C^0}
\big(c_{\psi_\infty}^{-(1+2s)/2}
\lfloor\psi_\infty'\rfloor_{s,2}+
c_{\psi_k}^{-(1+2s)/2}
\lfloor\psi_k'\rfloor_{s,2}\big),\notag
\end{align}
according
to \cite[Lemma A.6]{knappmann-etal_2023} and \cref{eq:erste-chain}, where the 
bilipschitz constants $c_{\psi_k}$ of the $\psi_k$ remain
uniformly positively bounded from below because of the convergence
$\psi_k\to\psi_\infty$ in $C^1$ as $k\to\infty$, and by
the fact that the
diffeomorphism $\psi_\infty$ has a strictly positive bilipschitz
constant $c_{\psi_\infty}$. Therefore, the supremum of 
 \cref{eq:upper-bound-second-seminorm} over $k\in\N$ is finite.

The $C^0$-norm in the first summand of the right-hand side
of \cref{eq:seminorm-basic-estimate} is bounded from above by 
$\frac{1}{c^2}$,
so
it remains to investigate the seminorm in that first summand:   
\begin{equation}
\label{eq:inverse_fracnorm}
\lfloor\psi_k'\circ\phi_k-\psi_\infty'\circ\phi_\infty\rfloor_{s,2}
\le\lfloor\psi_k'\circ\phi_k-\psi_\infty'\circ\phi_k\rfloor_{s,2}+
\lfloor\psi_\infty'\circ\phi_k-\psi_\infty'\circ\phi_\infty\rfloor_{s,2}.
\end{equation}
The first term in \cref{eq:inverse_fracnorm}
goes to zero as $k\to\infty$ since 
\begin{align*}
&\textstyle
\lfloor\psi'_{k}\circ\phi_k-\psi_\infty'\circ\phi_k\rfloor_{s,2}^2
=\iint_{[0,1]^2}\frac{|\psi'_{k}(\phi_k(x))-\psi_\infty'(\phi_k(x))
-\psi'_{k}(\phi_k(y))+\psi_\infty'(\phi_k(y))|^2}{\left|x-y \right|^{1+2s}}\d x\d y \\
&
\textstyle
= \iint_{[0,1]^2}\frac{|\psi'_{k}(v)-\psi_\infty'(v)-\psi'_{k}(w)
+\psi_\infty'(w)|^2}{\left|\psi_k(w)-\psi_k(v) \right|^{1+2s}} 
\psi'_{k}(w) \psi'_{k}(v)\d w\d v
\le c_{\psi_k}^{-(1+2s)}\|\psi'_k \|^2_{C^0}
\lfloor \psi_k'\!-\!\psi_\infty'\rfloor_{s,2}^2
\to 0,
\end{align*}
where we used the 
change of variables $w=\phi_k(x)$ and 
$v=\phi_k(y)$ and the fact that the bilipschitz constants $c_{\psi_k}$
are uniformly positively bounded from below.

For the square of the
second summand in \cref*{eq:inverse_fracnorm} we estimate 
the integrand of the corresponding double integral as
\begin{align*}
&
\textstyle
\frac{|\psi_\infty'(\phi_k(x))-\psi_\infty'(\phi_\infty(x))-(\psi_\infty'
(\phi_k(y))
-\psi_\infty'(\phi_\infty(y)))|^2}{\left|x-y \right|^{1+2s}}
\textstyle
\le \frac{|\psi_\infty'(\phi_k(x))-\psi_\infty'(\phi_k(y))|^2}{\left|x
-y \right|^{1+2s}}+\frac{|\psi_\infty'(\phi_\infty(x))
-\psi_\infty'(\phi_\infty(y))|^2}{\left|x-y \right|^{1+2s}}.
\end{align*}
The last summand is integrable since it equals the integrand of
the squared seminorm $\lfloor\psi_\infty'\circ\phi_\infty\rfloor_{s,2}$,
which can be estimated as 
 in \cref{eq:erste-chain},
and for the first term we have for $x\not = y$ (by setting $w=\phi_k(x)$ and $v=\phi_k(y)$)
\begin{align*}
\textstyle
\frac{|\psi_\infty'(\phi_k(x))-\psi_\infty'(\phi_k(y))|^2}{\left|x-y \right|^{1+2s}}
&
\textstyle
=\frac{|\psi_\infty'(w)-\psi_\infty'(v)|^2}{\left|\psi_k(w)-\psi_k(v) \right|^{1+2s}}
\le
c_{\psi_k}^{-(1+2s)}\frac{|\psi_\infty'(w)
-\psi_\infty'(v)|^2}{\left|w-v \right|^{1+2s}}
\le C\frac{|\psi_\infty'(w)
-\psi_\infty'(v)|^2}{\left|w-v \right|^{1+2s}},
\end{align*}
where we used again that the bilipschitz constants $c_{\psi_k}$ are 
uniformly positively bounded from below.
The last term is integrable since it is the integrand of
the squared seminorm $\lfloor\psi_\infty'\rfloor_{s,2}$.
So we found an integrable majorant for the integrand almost everywhere 
(as the set $\{(x,y)\in[0,1]^2 \mid x\not=y\}$ is a set of measure zero). 
By the dominated convergence theorem we have 
$\lfloor \psi_\infty'\circ\phi_k-\psi_\infty'\circ\phi_\infty
\rfloor_{s,2}\to 0$, and therefore by \cref{eq:seminorm-basic-estimate}
also $\lfloor\phi_k'-\phi_\infty'\rfloor_{s,2}\to 0$ as $k\to\infty$.
\end{proof}
Analogously to \cite[Corollary 1.17]{reiter_2009} we have the following corollary:
\begin{corollary}\label{cor:restriction-inverse-operator}
The restriction of the inversion operator $^{-1}$ to $H^{1+s}_\periodic((0,1))$ 
is a continuous map onto $H^{1+s}_\periodic((0,1))$.
\end{corollary}
Next we show that the composition operator is well-defined and continuous. 
\begin{lemma}
\label[lemma]{composition_continuous}
Let $s\in (\frac{1}{2},1]$. The composition operator 
\begin{align}\label{eq:composition-operator}
H_p^{1+s}((0,1))\times H^{1+s}(\R/\Z,\R^n)\to H^{1+s}(\R/\Z,\R^n),\quad
(\psi,h)\mapsto h\circ\psi,
\end{align}
is well-defined and continuous. 
\end{lemma}
\begin{proof}
For $s=1$ one finds a proof in \cite[Theorem 1.18]{reiter_2009}, so
as before we focus on $s\in (\frac{1}{2},1)$. The chain rule \cite[Lemma A.4]{knappmann-etal_2023} yields $h\circ\psi\in H^{1+s}(\R/\Z,\R^n)$, so that the composition operator \cref{eq:composition-operator}
is well-defined.

As to continuity, 
consider a sequence $(\psi_k,h_k)_k\subset H_p^{1+s}((0,1))
\times H^{1+s}(\R/\Z,\R^n)$ converging in $H^{1+s}\times H^{1+s}$ to a limiting 
pair $(\psi_\infty,h_\infty)\in H_p^{1+s}((0,1))\times H^{1+s}(\R/\Z,\R^n)$.
Then we have 
\begin{align*}
\textstyle
\|h_k\circ\psi_k-h_\infty\circ\psi_\infty\|_{C^0}
&\textstyle
\le\|h_k\circ\psi_k-h_\infty\circ\psi_k\|_{C^0}
+\|h_\infty\circ\psi_k-h_\infty\circ\psi_\infty\|_{C^0}\\
&\le \|h_k-h_\infty\|_{C^0}+\|h_\infty'\|_{C^0}
\|\psi_k-\psi_\infty\|_{C^0}\to 0,
\end{align*}
where we used, that $\psi_k$ is a diffeomorphism for each $k\in\N$. 
Similarly, for the  derivative,
\begin{align}
\|(h_k\circ\psi_k)'- (h_\infty\circ\psi_\infty)'\|_{C^0}&
\le\|\big(\|h_k'-h_\infty'\|_{C^0}
+\|h_\infty'\circ\psi_k-h_\infty'\circ\psi_\infty\|_{C^0} \big)\|\psi_k'\|_{C^0}
\notag\\
&+\|h_\infty'\circ\psi_\infty\|_{C^0}\|\psi_k'-\psi_\infty'\|_{C^0}\to 0\quad\As
k\to\infty,\label{eq:zwischen}
\end{align}
since
$h_\infty'$ is uniformly continuous on $\R/\Z$ and 
$\psi_k\to\psi_\infty$ in $C^1$. 
Lastly, 
the seminorm $\lfloor(h_k\circ\psi_k)'-(h_\infty\circ\psi_\infty)'\rfloor_{s,2}
$ is bounded from above by
\begin{align*}
&\lfloor(h_k'\circ\psi_k)\psi_k'-(h_\infty'\circ \psi_\infty)\psi_k'
\rfloor_{s,2}
+\lfloor(h_\infty'\circ\psi_\infty)\psi_k'-(h_\infty'\circ \psi_\infty)
\psi_\infty\rfloor_{s,2}\\
&\le
\big(\|h_k'-h_\infty'\|_{C^0}
-\|h_\infty'\circ\psi_k-h_\infty'\circ\psi_\infty\|_{C^0}\big)
\lfloor\psi_k'\rfloor_{s,2}\\
& + \big(
\lfloor h_k'\circ\psi_k-h_\infty'\circ\psi_k\rfloor_{s,2}+
\lfloor h_\infty'\circ\psi_k-h_\infty'\circ\psi_\infty\rfloor_{s,2}\big)
\|\psi_k'\|_{C^0}\\
& + \|h_\infty'\circ\psi_\infty\|_{C^0}\lfloor\psi_k'-\psi_\infty'\rfloor_{s,2}
+\lfloor h_\infty'\circ\psi_\infty\rfloor_{s,2}\|\psi_k'-\psi_\infty'\|_{C^0},
\end{align*}
by means of the product rule \cite[Proposition A.5]{knappmann-etal_2023}.
The $k$-dependent $C^0$-norms in the first and third line of the right-hand side
converge to zero as $k\to\infty$
as seen in 
\cref{eq:zwischen}, and the seminorms $\lfloor\psi_k'\rfloor_{s,2}$
 are uniformly bounded in $k$, as well as
$\lfloor h_\infty'\circ\psi_\infty\rfloor_{s,2}$ and
the norm $\|h_\infty'\circ\psi_\infty\|_{C^0}$.
Since also the norms $\|\psi_k'\|_{C^0}$ are uniformly bounded, it suffices to
show that the $k$-dependent seminorms in the middle line 
converge to zero. The first
seminorm can be treated as the first summand on the right-hand side
of \cref{eq:inverse_fracnorm} with the only difference that now the change
of variables reads as $w=\psi_k(x)$ and $v=\psi_k(y)$, so that one
has to use the bilipschitz constants $c_{\phi_k}=\|\psi_k'\|_{C^0}^{-1}$
of the inverse functions $\phi_k=\psi_k^{-1}$ satisfying
$c_{\phi_k}\ge
\frac{1}{2}\|\psi_\infty'\|_{C^0}>0$ for $k\gg 1$, because $\psi_k\to\psi_\infty$ in $C^1$. 
The second seminorm in that middle line resembles the second
summand of \cref{eq:inverse_fracnorm} and can be treated analogously 
using dominated
convergence, only now with the substitution $w=\psi_k(x)$ and $v=\psi_k(w)$ which
again requires the bilipschitz constants $c_{\phi_k}$ which are uniformly
positively bounded from below as $k\to\infty$. 
	\end{proof}
\begin{lemma}
\label[lemma]{lem:parametrization_continuous}
The operator 
\begin{equation}\label{eq:arclength-operator}
\textstyle
\varPsi \colon H_\reg^{1+s}(\R/\Z,\R^n)\to H^{1+s}_\periodic((0,1)),\quad
\gamma\mapsto \big( t\mapsto \frac{1}{\mathcal{L}(\gamma)} \int_{0}^{t}\left|\gamma'(\tau) \right|\d \tau \big),
\end{equation}
is well-defined and continuous.
	\end{lemma}
	\begin{proof}
It is clear that $\varPsi \, \g \in H^{1+s}_p((0,1))$ for any 
$\g\in H^{1+s}_\reg(\R/\Z,\R^n)$ so that $\varPsi$ is well-defined.
For a sequence of curves $(\gamma_k)_{k\in\N}\subset H_\reg^{1+s}(\R/\Z,\R^n)$ 
with $\gamma_k\to\gamma_\infty$ in $H_\reg^{1+s}(\R/\Z,\R^n)$ 
as $k\to\infty$, 
there exists a constant
 $c>0$, such that we have
 $\min\{|\gamma_k'(x)|,|\g_\infty'(x)|\}\ge c$ 
 for all $x\in\R/\Z$ and $k\in\N$, which implies $\LL(\g_k)\ge c$ for
 all $k\in\N$.
 In addition, $\psi_k \ceq \varPsi \, \gamma_k$ is of class  
 $C^1$ with $\psi_k'=\frac{1}{\mathcal{L}(\gamma_k)}\left|\gamma_k' \right|$
 for each $k\in\N$. 
Furthermore,  
$\mathcal{L}(\gamma_k)\to\mathcal{L}(\gamma_\infty)$ as $k\to\infty$
by the $C^1$-convergence $\g_k\to\g_\infty$; hence
\begin{align*}
&
\textstyle
\|\psi_k(\cdot)-\psi_\infty\|_{C^0}
= \norm[\Big]{ 
	\frac{\mathcal{L}(\gamma_\infty)\int_{0}^{\cdot}|\gamma_k'|\d\tau
	-\mathcal{L}(\gamma_k)\int_{0}^{\cdot}|\gamma_\infty'|\d\tau}{\mathcal{L}(
	\gamma_k)\mathcal{L}(\gamma_\infty)} 
}_{C^0}\\
&
\textstyle
\le\frac{1}{c}\int_{0}^{1}|\gamma_k'(\tau)-\gamma_\infty'(\tau)|\d\tau 
+ \frac{1}{c^2}| \mathcal{L}(\gamma_k)-\mathcal{L}(\gamma_\infty)|
\int_{0}^{1}|\gamma_\infty'(\tau)|\d\tau\to 0\quad\As k\to\infty.
\end{align*}
For the $C^0$-norm of the derivative we get
\begin{align*}
\textstyle
\|\psi_k'-\psi_\infty'\|_{C^0}&
\textstyle
=\big\|\frac{1}{\mathcal{L}(\gamma_k)}
|\gamma_k'| -\frac{1}{\mathcal{L}(\gamma_\infty)}|\gamma_\infty'|\big\|_{C^0}
\textstyle
\le\frac{1}{c}\|\gamma_k'-\gamma_\infty'\|_{C^0}
+\frac{1}{c^2}|\mathcal{L}(\gamma_k)-\mathcal{L}(\gamma_\infty)|
\|\gamma_\infty'\|_{C^0},
\end{align*}
which converges to zero as $k\to\infty$. Analogously,
\begin{align*}
\textstyle
\lfloor\psi_k'-\psi_\infty'\rfloor_{s,2}\le\frac{1}{c} 
\lfloor\gamma_k'-\gamma_\infty'\rfloor_{s,2}+\frac{1}{c^2}
|\mathcal{L}(\gamma_k)-\mathcal{L}(\gamma_\infty)|\lfloor\gamma_\infty'
\rfloor_{s,2}\to 0\quad\As k\to\infty. 
\end{align*}
	\end{proof}
\begin{refproof}[Proof of \Cref{thm:arclength}]
By 
\Cref{inversioniscontinuous,composition_continuous,lem:parametrization_continuous} 
the operator is continuous as the composition of continuous operators. 
Furthermore, for $\gamma\in H^{1+s}_\reg(\R/\Z,\R^n)$, 
\begin{align*}
\textstyle
|\ARC(\gamma)'(x)|&
\textstyle
=\big|\frac{1}{\mathcal{L}(\gamma)}
\left( \gamma\circ(\varPsi \, \gamma)^{-1}\right)'(x) \big|
=\frac{1}{\mathcal{L}(\gamma)}\left|\gamma'\left( \left( \varPsi \, \gamma \right)^{-1} (x)\right)(\left( \varPsi \, \gamma \right)^{-1})'(x) \right|\\
&
\textstyle
=\frac{1}{\mathcal{L}(\gamma)}\big|\frac{\gamma'\left( (\varPsi \, \gamma)^{-1}(x) \right)}{(\varPsi \, \gamma)'((\varPsi \, \gamma)^{-1}(x))} \big|=1.
\end{align*}
\end{refproof}

\begin{lemma}[Triple integrals]
\label{lem:triple-integrals}
Let $s\in (0,1)$, $\rho\in [1,\infty)$, and $\g\in W^{1+s,\rho}(\R/\Z,\R^n)$.
Then
\begin{equation}\label{eq:triple-integrals}
\int_{\R/\Z}\int_{-\frac{1}{2}}^\frac{1}{2}\int_0^1\frac{|\g'(u+\theta w)-\g'(u)|^\rho}{
|w|^{1+s\rho}}\d\theta\d w\d u\le \lfloor\g'\rfloor_{s,\rho}^\rho.
\end{equation}
\end{lemma}
\begin{proof}
Using Fubini and the change of variable $\tilde{w} \ceq \theta w$ with $\d \tilde{w}=
\theta\d w$ and $\tilde{w}(\pm\frac{1}{2})=\pm\frac{\theta}2$ one can bound 
the
triple integral in \cref{eq:triple-integrals} by means of $\theta\in [0,1]$
from above by
\begin{equation}\label{eq:triple-integral-upper-bound}\textstyle
\int_0^1\!\!\int_{\R/\Z}\int_{-\frac{1}{2}}^\frac{1}{2}\frac{|\g'(u+\tilde{w})-\g'(u)|^\rho}{|\tilde{w}|^{1+s\rho}}\d \tilde{w}\d u\d\theta\le 
\lfloor\g'\rfloor_{s,\rho}^\rho.
\end{equation}
\end{proof}





